% RMG-1
% Relational Medium Gravity
% Internal core draft, version 1
% SEARCH TAGS:
% RMG-CORE
% RELATIONAL-MEDIUM-GRAVITY
% BASE-MEDIUM
% PHANTOM-FREE
% ADM-GR-FIXED-POINT

\documentclass[11pt]{article}
\usepackage{amsmath,amssymb,amsthm}
\usepackage[margin=1in]{geometry}

\newtheorem{axiom}{Axiom}
\newtheorem{gate}{Gate}
\newtheorem{result}{Result}

\begin{document}

\begin{center}
{\Large RMG-1: Relational Medium Gravity}\\
\vspace{0.4em}
{\large რელაციური მედიუმის გრავიტაცია}\\
\vspace{0.4em}
Internal core draft, version 1\\
2026-05-10
\end{center}

\section*{Purpose}

RMG is the phantom-free rebuild of the previous ISPG line.

Core idea:
\[
\boxed{
\text{particles, rods, clocks and geometry are relational modes of one base-medium.}
}
\]

Gravity is not introduced as a fundamental phantom scalar field.
Instead:
\[
\boxed{
\text{GR/ADM is the low-energy fixed point of the base-medium's relational dynamics.}
}
\]

\section*{Logical status labels}

RMG-1 separates four kinds of statements:
\[
\boxed{\text{Axiom}}
\]
minimal starting property of the base-medium;
\[
\boxed{\text{Derived in RMG-1}}
\]
result obtained inside this file after accepting the stated assumptions;
\[
\boxed{\text{Target}}
\]
structure required by observation and used as the low-energy goal;
\[
\boxed{\text{Open theorem}}
\]
statement still to be proven from microscopic base-medium dynamics.

\section{Definition}

\[
\boxed{
\text{RMG = Relational Medium Gravity.}
}
\]

ქართული:
\[
\boxed{
\text{რელაციური მედიუმის გრავიტაცია.}
}
\]

Working definition:

RMG is a gravitational framework in which a single base-medium $\mathcal M$ gives rise to matter modes, clocks, rods, signals and operational geometry. The tested GR/ADM structure is treated as the low-energy fixed point of this medium, while the deeper task is to explain why that structure emerges.

\section{Core terms}

\[
\mathcal M
\]
is the base-medium. It is not assumed to be an ordinary scalar field.

\[
u
\]
is the compression/lapse potential. It is not a free phantom particle.

\[
s(x)
\]
is the universal local operational factor for clocks and rods.

\[
g_{\mu\nu}^{\rm eff}
\]
is the effective metric seen by low-energy matter modes.

\[
N,\quad N^i,\quad h_{ij}
\]
are the ADM lapse, shift and spatial metric variables in the low-energy fixed-point description.

\[
h_{ij}^{TT}
\]
is the propagating tensor/shear gravitational-wave sector.

\section{Minimal axioms}

\begin{axiom}[One base-medium]
All low-energy matter modes, rods, clocks, signals and geometric excitations are collective modes of one base-medium $\mathcal M$.
\end{axiom}

\begin{axiom}[Positive physical energy]
Physical excitations of $\mathcal M$ have Hamiltonian bounded below:
\[
H-H_{\rm vac}\ge0.
\]
There is no fundamental phantom degree of freedom.
\end{axiom}

\begin{axiom}[Operational relational invariance]
Coordinate labels are not observables. Only relations between physical modes are observables.
\end{axiom}

\begin{axiom}[Universal operational response]
In a slowly varying compression state, all stable low-energy matter modes share one local clock/rod factor:
\[
s=s(x).
\]
Thus all matter sees one effective metric.
\end{axiom}

\begin{axiom}[Square signal law]
Coordinate signal speed obeys:
\[
\frac{c_{\rm coord}}{c}=s^2.
\]
\end{axiom}

\begin{axiom}[Multiplicative layering]
If weak compression potentials add,
\[
u=u_1+u_2,
\]
then process factors multiply:
\[
A(u_1+u_2)=A(u_1)A(u_2).
\]
With weak-field normalization, this gives:
\[
A=e^{-u}.
\]
\end{axiom}

\begin{axiom}[Total Hamiltonian source]
The compression/Hamiltonian constraint is sourced by total Hamiltonian density:
\[
\mathcal H_{\rm total}.
\]
\end{axiom}

\begin{axiom}[No local preferred frame]
The base-medium may define a cosmological state, but no local low-energy experiment may detect an absolute medium rest frame beyond observational bounds.
\end{axiom}

\begin{axiom}[Constrained scalar/vector sectors]
The scalar compression variable and the vector flow variable are constraint/near-zone variables at low energy, not independent observable radiation modes.
\end{axiom}

\begin{axiom}[Tensor radiation]
The only propagating vacuum gravitational modes at low energy are transverse-traceless tensor modes:
\[
h_+,\qquad h_\times,
\]
with positive energy and propagation speed $c$.
\end{axiom}

\begin{axiom}[Constraint closure]
The low-energy constraints close under Poisson brackets:
\[
\{\mathcal C_A,\mathcal C_B\}
=
f_{AB}^{\;\;C}\mathcal C_C.
\]
The target realization is the ADM hypersurface deformation algebra.
\end{axiom}

\section{Core logical chain}

The intended chain is:
\[
\boxed{
\text{base-medium axioms}
\Rightarrow
\text{universal operational metric}
\Rightarrow
\text{ADM/GR fixed point}
\Rightarrow
\text{tested gravity}.
}
\]

RMG-1 explicitly checks only the weak-field middle part:
\[
\text{Hamiltonian source}
\Rightarrow
u=GM/(rc^2)
\Rightarrow
\gamma=\beta=1.
\]

The deepest step remains open:
\[
\boxed{
\text{base-medium axioms}
\Rightarrow
\text{ADM/GR fixed point}.
}
\]

\section{Static weak-field recovery}

Let one localized matter mode have positive total energy density:
\[
\rho_E(\mathbf x)\ge0.
\]

Its inertial mass is:
\[
M
=
\frac{1}{c^2}\int d^3x\,\rho_E(\mathbf x).
\]

RMG weak Hamiltonian constraint:
\[
\nabla^2u
=
-
\frac{4\pi G}{c^4}\rho_E.
\]

With $u(\infty)=0$:
\[
u(\mathbf x)
=
\frac{G}{c^4}
\int d^3x'\,
\frac{\rho_E(\mathbf x')}{|\mathbf x-\mathbf x'|}.
\]

Far field:
\[
\boxed{
u(r)=\frac{GM}{rc^2}+O(r^{-2}).
}
\]

Thus the gravitational charge is inertial mass because the source is total energy.

\section{Operational metric}

The weak-field operational metric target is:
\[
\boxed{
ds^2=-e^{-2u}c^2dt^2+e^{2u}\delta_{ij}dx^idx^j.
}
\]

Here:
\[
A=e^{-u},
\qquad
B=e^{u}.
\]

For $u\ll1$:
\[
g_{00}=-1+2u-2u^2+O(u^3),
\]
\[
g_{ij}=(1+2u+O(u^2))\delta_{ij}.
\]

Since:
\[
U=c^2u,
\]
we get:
\[
U(r)=\frac{GM}{r}.
\]

Therefore:
\[
\boxed{
\gamma=1,\qquad \beta=1.
}
\]

\section{Signal speed}

For radial null propagation:
\[
0=-e^{-2u}c^2dt^2+e^{2u}dr^2.
\]

So:
\[
\frac{c_{\rm coord}}{c}=e^{-2u}=A^2.
\]

This realizes:
\[
\boxed{
c_{\rm coord}\propto L_{\rm oper}^2.
}
\]

Status:
\[
\boxed{
\text{Derived in RMG-1 from the operational metric map.}
}
\]

\section{Linear dynamic split}

ADM variables:
\[
ds^2
=
-N^2c^2dt^2+h_{ij}(dx^i+N^idt)(dx^j+N^jdt).
\]

Weak split:
\[
N=1-u+O(u^2),
\]
\[
h_{ij}=(1+2u)\delta_{ij}+h_{ij}^{TT}+O(2),
\]
\[
N_i=\mathcal A_i+O(2).
\]

Scalar/compression constraint:
\[
\nabla^2u=-4\pi G\rho/c^2.
\]

Vector/flow constraint:
\[
\nabla^2\mathcal A_i
=
-
\kappa_VG\mathcal P_i^T/c^3.
\]

Tensor/shear radiation:
\[
\left(
\frac{1}{c^2}\partial_t^2-\nabla^2
\right)h_{ij}^{TT}
=
\kappa_TGT_{ij}^{TT}/c^4.
\]

Vacuum propagation:
\[
\left(
\frac{1}{c^2}\partial_t^2-\nabla^2
\right)h_{ij}^{TT}=0.
\]

Only:
\[
h_+,\qquad h_\times
\]
propagate in vacuum.

\section{Nonlinear completion}

The safe nonlinear completion is the Einstein-Hilbert/ADM fixed point:
\[
\boxed{
S_{\rm low}
=
\frac{c^4}{16\pi G}
\int d^4x\sqrt{-g}R
+
S_{\rm matter}[g,\Psi].
}
\]

ADM form:
\[
S_{\rm ADM}
=
\int dt\,d^3x
\left(
\pi^{ij}\dot h_{ij}
-
N\mathcal H
-
N^i\mathcal H_i
\right).
\]

Hamiltonian constraint:
\[
\mathcal H\approx0.
\]

Momentum constraints:
\[
\mathcal H_i\approx0.
\]

RMG interpretation:
\[
N=e^{-u}
\quad\text{is the local process/lapse factor},
\]
\[
N^i
\quad\text{is the flow/frame-dragging variable},
\]
\[
h_{ij}
\quad\text{is the operational spatial metric}.
\]

Status:
\[
\boxed{
\text{Target / safe fixed point, not yet derived from }\mathcal M.
}
\]

\section{What is new}

RMG does not try to replace GR in tested regimes.

It tries to explain:
\[
\boxed{
\text{why GR/ADM emerges from the base-medium.}
}
\]

The old scalar pressure idea survives only as:
\[
\boxed{
\text{compression/lapse interpretation of the Hamiltonian constraint.}
}
\]

It is not a free phantom scalar field.

\section{Statement status map}

\[
\begin{array}{c|c}
\text{Statement} & \text{Status} \\
\hline
\text{one base-medium }\mathcal M & \text{Axiom} \\
\text{positive physical energy} & \text{Axiom} \\
\text{universal operational factor }s & \text{Axiom / open microscopic derivation} \\
c_{\rm coord}/c=s^2 & \text{Axiom / observationally motivated} \\
A=e^{-u} & \text{Derived from multiplicative layering} \\
u(r)=GM/(rc^2) & \text{Derived from Hamiltonian source} \\
\gamma=1,\ \beta=1 & \text{Derived in weak field} \\
\text{only two tensor GW modes} & \text{Target axiom / observational gate} \\
\text{ADM/EH low-energy fixed point} & \text{Target / open theorem} \\
\text{derivation from }\mathcal M & \text{Main open theorem}
\end{array}
\]

\section{What is not yet proven}

RMG-1 is not a complete theory.

Still open:
\[
\text{derive ADM/Einstein-Hilbert from the base-medium axioms},
\]
\[
\text{derive universal matter coupling from microscopic modes},
\]
\[
\text{derive tensor/shear sector from medium mechanics},
\]
\[
\text{control preferred-frame effects},
\]
\[
\text{rebuild MOND/cosmology/particle sectors without phantom foundation}.
\]

\section{Pass/fail gates}

\begin{gate}[Core pass]
RMG-1 passes as a core framework if it can derive:
\[
S_{\rm EH}+S_{\rm matter}
\]
as the low-energy fixed point of $\mathcal M$, with positive physical Hamiltonian and only two vacuum tensor gravitational degrees of freedom.
\end{gate}

\begin{gate}[Core fail]
RMG-1 fails if the base-medium produces an observable preferred frame, species-dependent metric coupling, extra scalar/vector radiation, or a wrong-sign physical mode.
\end{gate}

\section*{Current status}

\[
\boxed{
\text{RMG-1 is a strong phantom-free core direction, not yet a closed theory.}
}
\]

Main theorem to prove:
\[
\boxed{
\text{base-medium relational axioms}
\Rightarrow
\text{ADM/Einstein-Hilbert low-energy fixed point}.
}
\]

\end{document}
